\subsection{A simultaneous fixed-point module obstruction}
\label{subsec:simultaneous-fixed-point}

\begin{theorem}
The fusion rings $R8\_3\_1203$, $R8\_3\_1350$, $R8\_3\_2842$, and
$R8\_3\_2885$ have no categorification over $\mathbb C$.
\end{theorem}
\begin{proof}
We give a general reduction and then its exact finite verification.
Let $g$ be an invertible object of order two fixing a simple object in a
categorification of a commutative fusion ring.  A fixed simple gives a
rank-one module category over $\operatorname{Vec}(C_2,\omega)$, so its
module associativity trivializes $[\omega]$.  Thus $A=\mathbf1\oplus g$
is a connected separable algebra.  In the category $\mathcal C_A$ of
right modules, a two-element $g$-orbit induces a single simple module,
whereas a fixed simple $V$ induces two simples $V_+,V_-$, each with
underlying object $V$.  These statements follow from induction
adjunction and semisimplicity.

The span of induced objects is invariant under every action matrix
$M_X$, and under its transpose $M_{X^*}$ by rigidity.  Its orthogonal
complement, with basis $V_+-V_-$ indexed by the fixed simples, therefore
carries matrices $S_X$ satisfying
\[
 S_{\mathbf1}=I,\qquad S_{X^*}=S_X^t,\qquad
 S_XS_Y=\sum_Z N_{XY}^Z S_Z.
\]
The corresponding two-by-two module block has both row and column sums
$a=N_{XV}^W$; hence it is
$\left(\begin{smallmatrix}u&a-u\\a-u&u\end{smallmatrix}\right)$.
Consequently $(S_X)_{VW}\in\{-a,-a+2,\ldots,a\}$.
The matrix $S_g$ is diagonal with entries $\pm1$.
If $gX=X$ and $N_{XV}^V$ is odd, then
$S_gS_X=S_X$ forces $(S_g)_{VV}=1$.  For the four rings in the theorem,
every fixed vertex admits such an $X$, and therefore $S_g=I$.

The fixed indices are $4,5,6,7$, with $g$ at index $3$ in the first two
rings and index $1$ in the last two.  The representation $S$ factors
through $R_{\mathbb Q}/(g-1)$, a six-dimensional rational algebra
generated by the self-dual basis object of index $4$.  The annihilating
polynomial for its image is
\[
\begin{cases}
(t^2-14t+4)(t^4+2t^3-6t^2-12t-4),&1203,2842,\\
(t^2-12t-14)(t^4-10t^2+12t-2),&1350,2885.
\end{cases}
\]
Every $S_X$ is a specified rational polynomial in $S_4$.  The exact
polynomial identities are checked independently in the original
rank-eight regular representation by multiplication by $I+N_g$.

We enumerate all symmetric integral $4\times4$ matrices $S_4$ having
the prescribed bounds and parities, test the displayed polynomial,
and reconstruct the other $S_X$.  A simultaneous bounded integral
representation never occurs:
\[
\begin{array}{c|r|r|r}
\text{ring}&\text{all matrices}&\text{polynomial survivors}&
\text{simultaneous survivors}\\\hline
1203&110592&32&0\\
1350&62208&16&0\\
2842&110592&32&0\\
2885&62208&16&0
\end{array}
\]
All computations use exact rational or unbounded integer arithmetic;
the finite bounds come from the module-block argument.  The discovery
enumerator uses a simultaneous diagonal-sign gauge and is replayed by
an independent enumerator without this reduction.  Both retain their
checks under optimized Python.  As a positive control, the known
model $R8\_3\_342=\operatorname{Rep}(C_{13}\rtimes C_6)$ passes every
simultaneous multiplication, duality, bound and parity check on its
two-dimensional difference representation.  The evidence and two
read-only verifiers are in
\texttt{evidence/rank8\_frontier\_20261021/quadratic/}.
\end{proof}

The argument uses only semisimple module-category theory (see Ostrik,
\emph{Module categories, weak Hopf algebras and modular invariants},
Transformation Groups \textbf{8} (2003), 177--206, and
Etingof--Gelaki--Nikshych--Ostrik, \emph{Tensor Categories}, Chapter~7).
In particular, neither pivotality nor unitarity is assumed.  The
conclusions for $1350$ and $2885$ strengthen the previous unitary-only
obstructions.

\begin{theorem}
The fusion ring $R8\_3\_921$ has no categorification over $\mathbb C$.
\end{theorem}
\begin{proof}
Use the same split module, now with $g$ at index $5$ and fixed indices
$3,4,6,7$.  Put $B=S_3$.  Since index $4$ is dual to $3$ and the ring
is commutative, $BB^t=B^tB$.  Thus $B$ is a real normal matrix.
All sixteen fixed-block coefficients of index $3$ are odd, and its four
diagonal coefficients are $1$.  Hence
$|\operatorname{tr}B|\leq4$ and $\|B\|_F^2\geq16$.
Its squarefree source annihilator is
\[
 p(t)=t(t^2-8t+6)(t^2+2t+4).
\]
The characteristic polynomial of this rational four-dimensional matrix
is $t^a(t^2-8t+6)^b(t^2+2t+4)^c$, where
$a+2b+2c=4$.  Its trace is $8b-2c$, and normality gives squared
Frobenius norm $52b+8c$.  The trace bound rules out $b>0$; the norm
bound then forces $c=2$.  Thus $q(B)=0$, where $q(t)=t^2+2t+4$.
The exact identity in the regular fusion representation
\[
 q(N_3)(-7N_1-7N_2-3N_3+30N_6-7N_7)=148N_6
\]
therefore forces $S_6=0$.  This contradicts the odd fixed-block entry
$N_{6,3}^{3}=1$.  The argument needs no condition $S_g=I$.
The verifier \texttt{normal\_split\_spectral.py} checks all source
identities and all six characteristic-factor possibilities.  Its
independent full-box search examines all $131072$ bounded matrices:
$2096$ are normal and $16$ satisfy $p(B)=0$, and every one of these
satisfies $q(B)=0$.  The matrix norm here is an ordinary exact
linear-algebra invariant; it introduces no categorical unitarity
assumption.
\end{proof}
